\newtoggle{withbonus} \togglefalse{withbonus} % set true to show bonus points in grade table

% setting underlying course name (Grundlagenfach Mathematik, §-Unterricht, \ldots)
\newcommand{\mycourse}{Grundlagenfach Informatik}
% name of the class
\newcommand{\myclass}{Klasse 3g}
% set date
\newcommand{\mydate}{10. Januar 2025}

\title{\textbf{Prüfung: Aus Daten lernen, endliche Automaten}}
\author{Thomas Graf\\
	\mycourse, \myclass}
\date{\mydate}
\maketitle
\thispagestyle{headandfoot}

\begin{center}
	\fbox{\parbox{140mm}{\centering
			\begin{itemize}
				\item Dauer der Prüfung: $45$ Minuten
				      %\item \textbf{Zeige in jeder Aufgabe unbedingt Deine Schritte!}
				\item Erlaubte Hilfsmittel:
				      \begin{itemize}
					      \item Schreibzeug
					      \item Taschenrechner
				      \end{itemize}
			\end{itemize}}}
\end{center}
\vspace{10mm}

\begin{center}
	Vorname: \rule{45mm}{0.8pt}\qquad Nachname: \rule{45mm}{0.8pt}
\end{center}
\vspace{20mm}

\begin{center}
	\iftoggle{withbonus} {
		\combinedgradetable[h][questions]
	} {
		\gradetable[h][questions]
	}
\end{center}
\vspace{15mm}
\begin{center}
	Note:
	\vspace{10mm}

	\rule{8mm}{1.2pt}
\end{center}
%\thispagestyle{empty}
\clearpage

\begin{questions}

	\section*{\textcolor{Green}{Theoretische Informatik}}
	\question
	\begin{parts}
		\part[1] Es seien $L_1 := \lrc{\lambda, a, aa}$ und $L_2 := \lrc{a, ba}$ bestimmen Sie die Anzahl der Elemente in der Menge $L_1L_2$.
		\vspace{3cm}
		\part[1]
		Geben Sie eine endliche Sprache $L$ an, sodass $\abs{L^2} < \abs{L}^2$.
		\vspace{3cm}
		\part[1]
		Ist die Beziehung $\lr{\lrc{a, b}^2}^*\subseteq \lr{\lrc{a, b}^*}^2$ wahr oder falsch? Begründen Sie Ihre Antwort.
        \droptotalpoints
	\end{parts}
	\clearpage

	\question{\textbf{Endlichen Automaten zeichnen 1}}

	Gegeben ist die Sprache
	\begin{align*}
		L_1 := \setcm{0x11y}{x,y\in\lrc{0,1}^*}.
	\end{align*}
	\begin{parts}
		\part[1] Nennen Sie ein kürzestes Wort in $L_1$.
		\begin{solution}
			Das Wort $011$ ist \textbf{das} (eindeutige) kürzeste Wort in $L_1$.
		\end{solution}
		\vspace{5cm}
		\begin{solution}
			\begin{figure}[H]
				\centering
				\begin{tikzpicture}[->,>=stealth',shorten >=1pt,node distance=3.2cm,semithick]
					\tikzstyle{every state}=[fill=green!10,draw=black,text=black,minimum size=1.4cm]

					\node[initial,state] (q0) {$q_{0}$};
					\node[state] (q1) [above of=q0] {$q_{1}$};
					\node[state] (q2) [right of=q0] {$q_{2}$};
					\node[state] (q3) [right of=q2] {$q_{3}$};
					\node[state,accepting] (q4) [right of=q3] {$q_{4}$};

					\path (q0) edge [left] node {$1$} (q1)
					(q0) edge [below] node {$0$} (q2)

					(q1) edge [loop above] node {$0,1$} (q1)

					(q2) edge [below] node {$1$} (q3)
					(q2) edge [loop above] node {$0$} (q2)

					(q3) edge [above, bend right] node {$0$} (q2)
					(q3) edge [below] node {$1$} (q4)

					(q4) edge [loop above] node {$0,1$} (q4);
				\end{tikzpicture}
				\caption{Endlicher Automat (in Diagrammform) für die Sprache $L_1$}
				\label{fig:L1}
			\end{figure}
		\end{solution}
		\part[2] Entwerfen Sie einen endlichen Automaten (in Diagrammform), welcher $L_1$ akzeptiert.
	\end{parts}
	\droptotalpoints

	\clearpage \question{\textbf{Endlichen Automaten zeichnen 2}}

	Gegeben ist die Sprache \begin{align*} L_2 := \setcm{x1y}{x,y\in\lrc{0,1}^*, \abs{y} = 1}.
	\end{align*}
	\begin{parts}
		\part[1] Nennen Sie ein kürzestes Wort in $L_2$.
		\begin{solution}
			Die beiden Wörter $10$ und $11$ sind die beiden kürzesten Wörter in $L_2$.
		\end{solution}
		\vspace{5cm}
		\part[3] Entwerfen Sie einen endlichen Automaten (in Diagrammform), welcher $L_2$ akzeptiert. \textcolor{Red}{Achtung}: Beachten Sie, dass $110\in L$ aber $100\notin L$.
		\begin{solution}
			\begin{figure}[H]
				\centering
				\begin{tikzpicture}[->,>=stealth',shorten >=1pt,node distance=5.2cm,semithick]
					\tikzstyle{every state}=[fill=green!10,draw=black,text=black,minimum size=1.4cm]

					\node[initial,state] (q0) {$q_{0}$};
					\node[state] (q1) [right of=q0] {$q_{1}$};
					\node[state,accepting] (q2) [right of=q1] {$q_{2}$};
					\node[state,accepting] (q3) [below right of=q1] {$q_{3}$};

					\path (q0) edge [loop below] node {$0$} (q0)
					(q0) edge [above] node {$1$} (q1)

					(q1) edge [below, bend right] node {$0$} (q2)
					(q1) edge [below] node {$1$} (q3)

					(q2) edge [above, bend right] node {$0$} (q0)
					(q2) edge [above] node {$1$} (q1)

					(q3) edge [right] node {$0$} (q2)
					(q3) edge [loop below] node {$1$} (q3);
				\end{tikzpicture}
				\caption{Endlicher Automat (in Diagrammform) für die Sprache $L_2$}
				\label{fig:L2}
			\end{figure}
		\end{solution}
	\end{parts}
	\droptotalpoints
	\clearpage

	\question[2]{\textbf{Anzahl der Elemente in einer Sprache bestimmen}}

	Es sei $\Sigma = \lrc{0, 1, 7}$ ein Alphabet und $n$ eine gegebene natürliche Zahl. Bestimmen Sie die Anzahl der verschiedenen Wörter der Länge $n$ über $\Sigma$, welche das Teilwort $7$ enthalten.
	\begin{solution}
		Insgesamt gibt es genau $\abs{\Sigma}^n = 3^n$ Wörter über $\Sigma$. Davon enthalten genau $2^n$ Wörter nicht $7$ als Teilwort. Die gesuchte Anzahl ist somit $3^n - 2^n$.
	\end{solution}
	\droptotalpoints
	\clearpage

	\question{\textbf{\red{Nicht reguläre} Sprache}}

	Gegeben sei die Sprache
	\begin{align*}
		L_3 := \setcm{w\$\$\$w^2}{w\in\lrc{0, 1}^*}
	\end{align*}
	über dem Alphabet $\lrc{0, 1, \$}$.

	(Zur Vermeidung von Missverständnissen: Für beispielsweise $w = 011$ ist $w^2 = (011)^2 = 011011$.)

	\begin{parts}
		\part[1] Nenne genau drei verschiedene Wörter, die in $L_3$ liegen.
		\begin{solution}
			\begin{align*}
				aaabaa, baabab, ab
			\end{align*}
		\end{solution}
		\vspace{3cm}
		\part[3] Zeige, dass $L_3$ \red{nicht regulär} ist.

		\begin{solution}
			Angenommen, $L_3$ sei regulär. Dann existiert ein endlicher Automat
			\begin{align*}
				A = (Q, {a,b}, \delta, q_0, F)
			\end{align*}
			und $A$ akzeptiert $L_3$. Sei jetzt $m := \abs{Q}$, also die Anzahl der verschiedenen Zustände des endlichen Automaten
			$A$. Wir betrachten die $m+1$ verschiedenen Wörter
			\begin{align*}
				a^1ab, a^2ab, \ldots, a^{m+1}ab.
			\end{align*}
			Dies sind mehr Worte (genau ein Wort mehr), als $A$ Zustände hat. Nach dem Schubfachprinzip (pigeonhole principle) existieren
			$i,j\in\lrc{1,2,\ldots,m+1}$ mit $i<j$, sodass
			\begin{align*}
				\hat{\delta}\lr{q_0,a^i} = \hat{\delta}\lr{q_0,a^j}.
			\end{align*}
			Nach der im Unterricht bewiesenen Aussage gilt
			\begin{align*}
				a^iz\in L_3 \iff a^jz\in L_3
			\end{align*}
			für alle $z\in\lrc{a,b}^*$. Doch dies gilt nicht, da zum Beispiel für die Wahl $z:=aba^i$
			\begin{align*}
				a^iz = a^iaba^i\in L_3 \quad\text{aber}\quad a^jz = a^jaba^i\notin L_3
			\end{align*}
			gilt, da $a^i$ zwar die Umkehrung von $a^i$ ist aber nicht von $a^j$. Damit war unsere ursprüngliche Annahme falsch und es existiert kein
			endlicher Automat für die Sprache $L_3$.
		\end{solution}
	\end{parts}
	\droptotalpoints

	\clearpage
	\section*{\textcolor{Green}{Beschreibende Statistik}}
	\question
	\begin{parts}
		\part[1] Nennen Sie ein Beispiel zweier Merkmale aus dem Alltag, die stark negativ korreliert sind.
		\begin{solution}
			Im Unterricht haben wir das Beispiel der starken positiven Korrelation zwischen Körpermasse und Körpergrösse gesehen.
		\end{solution}
		\vspace{4cm}
		\part[1] Nennen Sie ein Beispiel zweier Merkmale aus dem Alltag, die stark positiv korreliert sind.
		\vspace{4cm}
		\begin{solution}
			Zum Beispiel sind Aussentemperatur und Heizkosten stark negativ korreliert.
		\end{solution}
		\part[1] Erläutern Sie den Satz
		\begin{center}
			\textit{Korrelation impliziert nicht Kausalität.}
		\end{center}
		mithilfe eines konkreten Beispiels aus dem Alltag.
		\begin{solution}
			\url{https://www.tylervigen.com/spurious-correlations}
		\end{solution}

		\droptotalpoints
	\end{parts}

	\clearpage
	\question{\textbf{Punktklassifikator}}

	Wir haben die Längen (in Metern) von 5 Exemplaren der Pflanzenart \textit{Gefleckter Schierling} und von 4 Exemplaren der Pflanzenart \textit{Echter Kerbel} bestimmt:
	\begin{align*}
		\text{Längen der Gefleckten Schierlinge:}\quad & (3.04, 2.93, 2.94, 3.20, 3.17) \\
		\text{Längen der Echten Kerbel:}\quad          & (4.41, 4.30, 4.33, 4.38).
	\end{align*}

	\begin{parts}
		\part[1]
		Bestimmen Sie den Punktklassifikator für diese Situation.
        \vspace{4cm}
		\part[1]
		Wir führen eine weitere Messung der Länge eines neuen Exemplars eines Echten Kerbels durch. Sei $x$ diese weitere gemessene Länge. Welche Bedingung muss $x$ erfüllen, damit der Punktklassifikator für diese neue Situation nicht existieren würde?

        \droptotalpoints
	\end{parts}

\end{questions}
\clearpage
(leere Seite für Notizen)
\clearpage
(leere Seite für Notizen)
